\section{Perspectiva de ingeniería: por qué PPO se convirtió en la opción predeterminada}
\subsection{El compromiso entre complejidad, estabilidad y viabilidad de implementación}
Aunque el curso dedica solo 15 minutos a explicar la evolución de PG, TRPO y PPO, la conclusión de ingeniería que hay detrás es muy clara: PPO se volvió popular no porque sea absolutamente óptimo en todos los benchmarks, sino porque logra el compromiso más equilibrado entre \textbf{estabilidad, umbral de implementación y eficiencia de entrenamiento}.

\begin{table}[H]
\centering
\begin{tabular}{p{2.6cm}p{3.5cm}p{3.7cm}p{3.2cm}}
\toprule
Método & Ventaja & Costo principal & Cuándo usarlo \\
\midrule
PG & El más directo conceptualmente & Varianza alta, actualizaciones inestables & Para entender el principio básico \\
TRPO & Restricción teórica más fuerte & Aproximación de segundo orden e implementación compleja & Más adecuado para derivaciones de investigación \\
PPO & Estable y simple de implementar & Aún requiere ajustar el clip y la ventaja con precisión & Punto de partida predeterminado en la práctica de ingeniería \\
\bottomrule
\end{tabular}
\caption{Compromisos típicos de PG, TRPO y PPO en la ingeniería}
\end{table}

\begin{warningbox}{No tomes a PPO como un «algoritmo libre de ajuste de hiperparámetros»}
PPO solo sustituye las restricciones complejas de TRPO por un objetivo de recorte más fácil de usar, lo cual no significa que el entrenamiento se vuelva estable de manera automática. El tamaño del lote, la normalización de la ventaja, la escala del reward y la calidad del rollout siguen afectando directamente los resultados.
\end{warningbox}

\subsection{Resumen de la sección}
Desde la perspectiva de la puesta en producción de ingeniería, PPO es una línea base de alta relación costo-beneficio: primero se usa para armar el ciclo cerrado de entrenamiento, y después se decide si es necesario introducir restricciones o reglas de actualización más complejas; esto suele ser más realista que buscar desde el inicio el «algoritmo más fuerte».

\newpage
